%%%%%%%%%%%%%%%%%%%%%%%%%%%%%%%%%%%%%%%%%%%%%%%%%%%%%%%%%%%%
%% CHAPTER 1 -- STATISTICAL LEARNING FOUNDATIONS FOR LLMs
%%%%%%%%%%%%%%%%%%%%%%%%%%%%%%%%%%%%%%%%%%%%%%%%%%%%%%%%%%%%

\chapter{Statistical Learning Foundations for Large Language Models}
\label{sec:ch1}

\section*{Setting the Scene}
\addcontentsline{toc}{section}{Setting the Scene}

\paragraph{The problem.} Consider for a moment what it means to
predict the next word. You are reading this sentence. Your eye has
not yet reached its end, and yet some part of you is already
guessing how it will finish. The guess is not certain. It is not
arbitrary either. It is shaped by everything you have ever read, and
by the few words that came just before. If we wanted to teach a
machine to do the same, we would need to give it a way to assign,
to each possible continuation, a number that captures \emph{how
likely} that continuation is given everything it has seen so far.
That is the entire problem of language modelling, written in a
single sentence. The remainder of this chapter, and the chapters
that follow, are the elaboration.

\paragraph{How we got here.} The idea that one symbol in a sequence
depends on the symbols before it is older than the computer. In
1913 the Russian mathematician Andrey Markov sat down with a copy
of Pushkin's \emph{Eugene Onegin} and counted the alternation of
vowels and consonants. He counted twenty thousand letters by
hand. He did this, in part, to make a point against a colleague
who had argued that the law of large numbers required independent
trials; Markov wanted to show that one could relax independence
into a chain of conditional dependence and still do useful
mathematics. He gave us a model of language as a sequence of
dependencies a full century before we had the data or the hardware
to make that model interesting.

Claude Shannon, in 1948, wrote the paper that founded information
theory. He used English text as one of his running examples.
Shannon estimated the entropy of English at roughly one bit per
letter --- that is, on average each letter of an English string
carries about one bit of new information given everything before
it. He observed that a model that captures more of this structure
will, by definition, assign higher probability to natural text. He
gave us a yardstick.

For half a century after Shannon the field crawled. The n-gram
era of speech recognition counted word triples on hard drives
that filled rooms. Frederick Jelinek, working at IBM, gave us the
phrase ``every time I fire a linguist my speech recognizer's
performance improves'' --- a statement about the willingness of
counts to do the work that elaborate theory had failed to do.
Yoshua Bengio and his coauthors, in 2003, built a neural network
that learned a continuous representation of words and used it to
predict the next one. Almost nobody outside the field noticed.
Mikolov and his colleagues, in 2013, simplified the idea into
word2vec. The geometry of the resulting word vectors --- in which
``king minus man plus woman'' lands near ``queen'' --- caught
public attention in a way that the underlying technology, by then
ten years old, had not. The next decade and a half compressed:
sequence-to-sequence translation in 2014, attention in 2015, the
Transformer in 2017, and then the GPT line. By the end of 2022
the kind of model that had been a research curiosity for fifty
years was a consumer product.

\paragraph{Where we stand.} The lineage converged on a remarkably
simple recipe. Take a neural network. Show it a sequence of
tokens. Ask it to predict, for each position in the sequence, a
probability distribution over the next token. Compare its
prediction to the token that actually came next. Penalise it in
proportion to how surprised it was by the truth. That penalty,
averaged over a corpus, is called the \emph{cross-entropy} loss.
The mathematics in the rest of this chapter unpacks why this
particular penalty is the right one. The short version is that
cross-entropy is what you always end up at when you ask, formally,
``minimise the surprise of seeing the data we actually saw, and
nothing more.'' It is a maximum-likelihood estimator. It is a
proper scoring rule. It is the information-theoretic divergence
between two distributions (KL divergence is asymmetric, hence not a
true distance), plus a constant the model cannot control. We will see all three of these views in the pages that
follow, because each one buys a different intuition.

\paragraph{What goes wrong.} A model trained on cross-entropy is,
by construction, the most plausible thing the model can say. The
most plausible thing is not the same as the true thing. It is not
the same as the most useful thing. It is not the same as the
thing the user wanted, or the thing the user would have wanted if
they had thought about it longer. It is, narrowly, the
continuation that the training corpus would on average have
produced. We give a polite name to the gap between this and the
truth: we call it \emph{hallucination}. It is not a bug. It is the
loss surface doing exactly what we asked of it.

There are several specific consequences of this gap. The first is
that the model has no native, externally calibrated task-level
confidence in the sense a human reader would mean. It has a
probability distribution over
next tokens, which is not the same thing. We will distinguish
\emph{probability} (the model's number), \emph{confidence} (how
strongly the model commits to that number), and
\emph{correctness} (whether the chosen token is in fact right).
The three are routinely conflated by users, and that conflation is
the source of more downstream errors than any architectural
choice. The second consequence is that the model is what we call
\emph{calibrated}, or, more often, \emph{miscalibrated}: its
confidence levels do not reliably correspond to its accuracy
rates. We can measure this. We can sometimes correct it. We
cannot eliminate it.

The mathematics of the next sections is the bookkeeping for what
we have just described. We will derive cross-entropy from
maximum likelihood, identify it with information-theoretic
divergence, work an example with a softmax, define
calibration error formally, and watch how each of these objects
is one face of the same construction. By the end of this chapter
the reader who began with the intuition that an LLM ``just
predicts the next word'' will have a sturdier intuition: an LLM
encodes a conditional distribution, and reading its outputs
through that lens explains, in advance, most of the surprises the
rest of the book will document.

\section{Learning Objectives}

By the end of this chapter, you should be able to:
\begin{itemize}
  \item Formally interpret a Large Language Model (LLM) as a conditional probabilistic model over sequences of discrete symbols.
  \item Re-derive the cross-entropy loss used in next-token prediction from first principles, both as a maximum-likelihood estimator and as a consequence of the information-theoretic identity $H(p,q) = H(p) + \KL(p \parallel q)$.
  \item Work through a softmax computation numerically and show why the logit representation is invariant under additive constant shifts (justifying the log-sum-exp trick).
  \item Distinguish the three quantities that are routinely conflated in downstream discussions: \emph{probability}, \emph{confidence}, and \emph{correctness}.
  \item Explain why cross-entropy minimization does not, on its own, imply semantic correctness, and why this is the root cause of several characteristic LLM failure modes.
  \item Reason quantitatively about model calibration using reliability diagrams, Expected Calibration Error (ECE), and temperature scaling.
  \item Explain the information-theoretic interpretation of training objectives, including the equivalence between cross-entropy and perplexity.
  \item State the standard proper scoring rules used in probabilistic forecasting and explain why log-loss is a natural choice.
  \item Identify the human-factor risks that arise when probabilistic outputs are interpreted as authoritative assertions, and describe first-line mitigations.
\end{itemize}

\section{LLM Components Covered}

This chapter focuses on the statistical substrate underlying all LLM systems. Every concept introduced here will reappear in later chapters under increasingly engineering-heavy disguises, so the aim is to establish vocabulary and intuition that later chapters can take as given.

\textbf{Modeling and Training}
\begin{itemize}
  \item Conditional probability modeling over discrete tokens.
  \item Maximum likelihood estimation (MLE) as a training principle.
  \item Cross-entropy as the empirical-risk objective.
  \item Softmax parameterization and logit representation.
  \item Calibration and uncertainty quantification.
\end{itemize}

\textbf{Systems Perspective}
\begin{itemize}
  \item Interpretation of model outputs as distributions rather than assertions.
  \item Implications of probabilistic behavior for downstream system design, monitoring, and human-in-the-loop review.
\end{itemize}

\section{Notation and Assumptions}
\label{sec:ch1_notation}

\begin{notationbox}
We fix a finite vocabulary $V$ of size $|V|$. Sequences are written
$x = (x_1, \dots, x_T)$ with $x_t \in V$ and length $T \in \mathbb{N}$.
Random variables are uppercase (e.g.\ $X_t$); realisations are lowercase
(e.g.\ $x_t$). The shorthand $x_{<t} := (x_1, \dots, x_{t-1})$ denotes
the prefix.

Model parameters are $\theta \in \mathbb{R}^P$ with $P$ the total
parameter count; the conditional the model realises is $q_\theta(\cdot
\mid x_{<t}) := P_\theta(X_t = \cdot \mid X_{<t} = x_{<t})$. The dataset
is $\mathcal{D} = \{x^{(i)}\}_{i=1}^{N}$ with sequence lengths $T_i$ and
total tokens $N_{\mathrm{tok}} = \sum_i T_i$. The empirical measure
$\hat{p}$ is the uniform distribution on the set of observed
(context, token) pairs $\{(x^{(i)}_{<t}, x^{(i)}_t)\}$.

Expectations $\E_p[\cdot]$ are with respect to the distribution written
in the subscript; unsubscripted $\E$ uses the empirical measure
$\hat{p}$ by default. Logarithms are natural; entropy and cross-entropy
are in $\nats$ unless marked $\log_2$, in which case the unit is
$\bits$.

\textbf{Assumptions used in this chapter.}
\begin{itemize}
  \item[A1.1a] Sequences in $\mathcal{D}$ are i.i.d.\ realisations of an
  unknown sequence distribution. (Sequence-level independence.)
  \item[A1.1b] The token-level empirical loss treats the union of
  (context, token) pairs across sequences as i.i.d.\ under the
  empirical measure $\hat{p}$. This is a derived computational
  convenience: token-level pairs within a sequence are not
  independent under the true sequence distribution, but their
  contributions to the average NLL are unweighted by sequence
  identity.
  \item[A1.2] The vocabulary $V$ is finite and fixed; distributions over
  $V$ are strictly positive after softmax (so $\log q_\theta$ is finite
  everywhere).
  \item[A1.3] Calibration is evaluated on the same conditioning
  distribution as training (no distribution shift between the calibration
  split and the operating distribution).
\end{itemize}
\end{notationbox}

\section{Conceptual Core: What Kind of Object Is an LLM?}

\Def\ A Large Language Model defines a probability distribution over
sequences of discrete symbols drawn from the finite vocabulary $V$.
The object of study is the joint distribution
\begin{equation}
P_\theta(x_1, x_2, \ldots, x_T),
\label{eq:joint}
\end{equation}
parameterised by $\theta \in \mathbb{R}^P$. For modern subword
vocabularies $|V|$ is typically on the order of $10^4$ to $10^5$.

\Ident\ The sample space for (\ref{eq:joint}) has cardinality $|V|^T$,
which is combinatorially infeasible even for modest $T$. The chain rule
of probability factors (\ref{eq:joint}) into an ordered product of
conditionals:
\begin{equation}
P_\theta(x_1, \ldots, x_T) = \prod_{t=1}^{T} P_\theta(x_t \mid x_1, \ldots, x_{t-1}).
\label{eq:chain}
\end{equation}
This identity is exact and holds for any distribution over sequences.

\Cons\ An LLM commits to a particular parameterisation of the
conditional $P_\theta(x_t \mid x_{<t})$ (the neural architecture,
Chapter~5), and the weights $\theta$ are learned by fitting this
conditional to empirical text (Chapter~6). Everything else---the
reasoning-like behaviour, the hallucinations, the striking context
sensitivity---is a downstream consequence of this one modelling choice.

\Emp\ An LLM does \emph{not} explicitly store facts, rules, or
symbols. It stores numerical parameters that, together with its
architecture, implement a family of conditional distributions. The
model's ``knowledge'' is whatever makes those conditionals assign high
probability to corpus text. Many behaviours that look like reasoning
failures---confident fabrications, brittle chains of thought,
inconsistency under paraphrase---follow directly from reading the
model's outputs through a probabilistic lens rather than a symbolic
one.

\subsection{Why Left-to-Right?}

The chain rule in equation (\ref{eq:chain}) is written left-to-right, but it could just as well be written right-to-left, or in any other ordering. Any particular factorization is mathematically equivalent on the joint. Neural language models nonetheless commit to a canonical ordering---left-to-right for GPT-style autoregressive models, or a masked variant for BERT-style bidirectional encoders---because the \emph{parameter sharing} induced by that ordering is what makes the model tractable. We return to this choice in Chapter~6, where we contrast autoregressive, masked, and prefix-LM objectives.

\section{Maximum Likelihood and Next-Token Prediction}

\subsection{Training Objective}

\Def\ The \emph{maximum-likelihood estimator} (MLE) under model
$P_\theta$ and dataset $\mathcal{D}$ is
$\hat\theta := \arg\max_\theta P_\theta(\mathcal{D})$. The
\emph{token-normalised negative log-likelihood} (NLL) is
$\mathcal{L}_{\mathrm{NLL}}(\theta) := -\frac{1}{N_{\mathrm{tok}}}
\log P_\theta(\mathcal{D})$.

\begin{derivation}[: MLE to token-normalised NLL]
Using A1.1a (i.i.d.\ sequences) and the chain rule (\ref{eq:chain}):
\begin{align}
P_\theta(\mathcal{D})
  &= \prod_{i=1}^{N} P_\theta(x^{(i)})
  && \text{(A1.1a: i.i.d.\ sequences)} \label{eq:mle-1} \\
  &= \prod_{i=1}^{N} \prod_{t=1}^{T_i}
       P_\theta\!\big(x^{(i)}_t \mid x^{(i)}_{<t}\big)
  && \text{(chain rule, \ref{eq:chain})} \label{eq:mle-2} \\
  \log P_\theta(\mathcal{D})
  &= \sum_{i=1}^{N} \sum_{t=1}^{T_i}
       \log P_\theta\!\big(x^{(i)}_t \mid x^{(i)}_{<t}\big)
  && \text{($\log$ is a monotone homomorphism)} \label{eq:mle-3} \\
  \hat\theta
  &= \arg\max_\theta
       \sum_{i=1}^{N} \sum_{t=1}^{T_i}
       \log P_\theta\!\big(x^{(i)}_t \mid x^{(i)}_{<t}\big)
  && \text{(definition of $\hat\theta$)} \label{eq:mle-4} \\
  &= \arg\min_\theta
       \underbrace{\frac{-1}{N_{\mathrm{tok}}}
       \sum_{i=1}^{N} \sum_{t=1}^{T_i}
       \log P_\theta\!\big(x^{(i)}_t \mid x^{(i)}_{<t}\big)}_{=\;\mathcal{L}_{\mathrm{NLL}}(\theta)}
  && \text{(divide by $-N_{\mathrm{tok}}$)} \label{eq:mle}
\end{align}
Division by $-N_{\mathrm{tok}}$ is a monotone transformation, so the
$\arg\max$ of the log-likelihood equals the $\arg\min$ of the NLL.
Normalisation by $N_{\mathrm{tok}}$ gives a per-token quantity that is
comparable across datasets and batch sizes.
\end{derivation}

\Cons\ Equation (\ref{eq:mle}) is the training objective for
essentially every modern language model. Every innovation in the last
decade---Transformers, sparse mixtures, alignment, retrieval---is
layered on top of this one expression.

\subsection{Equivalence with Cross-Entropy}

\Def\ Let $\hat p$ denote the empirical measure on (context, token)
pairs induced by $\mathcal{D}$ (a uniform distribution on the
$N_{\mathrm{tok}}$ observed pairs). For two distributions $p, q$ on the
same space, the \emph{cross-entropy} is
$H(p, q) := \E_{x \sim p}[-\log q(x)] = -\sum_x p(x) \log q(x)$.

\begin{derivation}[: NLL is empirical cross-entropy]
\begin{align}
\mathcal{L}_{\mathrm{NLL}}(\theta)
  &= -\frac{1}{N_{\mathrm{tok}}} \sum_{i,t}
       \log q_\theta\!\big(x^{(i)}_t \mid x^{(i)}_{<t}\big)
  && \text{(from \ref{eq:mle})} \\
  &= -\E_{(x_{<t}, x_t) \sim \hat p}\!\left[
       \log q_\theta(x_t \mid x_{<t}) \right]
  && \text{(uniform average $=$ $\hat p$-expectation)} \\
  &= H(\hat p,\, q_\theta)
  && \text{(definition of cross-entropy)}.
\label{eq:nll-ce}
\end{align}
\end{derivation}

\Ident\ Minimising NLL is identical to minimising cross-entropy against
$\hat p$. This is a general fact about MLE for any probabilistic
model whose likelihood factorises over the data --- it does not
depend on the parameterisation belonging to an exponential family,
and it is not specific to language modelling.

\paragraph{Connection to KL divergence.}
\Ident\ For any distributions $p, q$ on a common finite support with
$q > 0$,
\begin{equation}
H(p, q) = H(p) + \KL(p \parallel q),
\label{eq:ce-kl}
\end{equation}
where $\KL(p \parallel q) := \sum_x p(x) \log \frac{p(x)}{q(x)}$.

\begin{derivation}[: $H(p,q) = H(p) + \KL(p \parallel q)$]
\begin{align}
H(p, q)
  &= -\sum_x p(x) \log q(x)
  && \text{(definition of cross-entropy)} \\
  &= -\sum_x p(x) \log p(x)
       + \sum_x p(x) \log \frac{p(x)}{q(x)}
  && \text{(add and subtract $\sum_x p \log p$)} \\
  &= H(p) + \KL(p \parallel q)
  && \text{(definitions of $H$ and $\KL$)}.
\end{align}
\end{derivation}

\Cons\ Since $H(p)$ does not depend on $\theta$, minimising
$H(p, q_\theta)$ over $\theta$ is equivalent to minimising
$\KL(p \parallel q_\theta)$. In training, $p = \hat p$, so we are
driving $q_\theta$ toward the empirical distribution in KL divergence.

\Emp\ The target of cross-entropy training is the \emph{data}
distribution, not the \emph{truth}. If the corpus contains systematic
errors or biases, a low-cross-entropy model faithfully reproduces them.
This double edge reappears in every alignment and governance
discussion downstream.

\subsection{Softmax Parameterization and Numerical Stability}

\Def\ The conditional $q_\theta(\cdot \mid x_{<t})$ is realised by the
model as a vector of real-valued \emph{logits}
$z \in \mathbb{R}^{|V|}$ computed by the final layer, followed by a
softmax:
\begin{equation}
q_\theta(x_t = v \mid x_{<t})
  = \frac{\exp(z_v)}{\sum_{v' \in V} \exp(z_{v'})}.
\label{eq:softmax}
\end{equation}
A1.2 (strictly positive softmax outputs) holds automatically because
$\exp(z_v) > 0$.

\paragraph{Shift invariance.}
\Ident\ For every $c \in \mathbb{R}$, shifting every logit by $c$
leaves the softmax unchanged:
$\mathrm{softmax}(z + c\,\mathbf{1}) = \mathrm{softmax}(z)$.

\begin{derivation}[: softmax shift invariance]
\begin{align}
\mathrm{softmax}(z + c\,\mathbf{1})_v
  &= \frac{\exp(z_v + c)}{\sum_{v'} \exp(z_{v'} + c)}
  && \text{(definition, with shifted argument)} \\
  &= \frac{\exp(c)\,\exp(z_v)}{\exp(c)\,\sum_{v'} \exp(z_{v'})}
  && \text{($\exp$ is multiplicative)} \\
  &= \frac{\exp(z_v)}{\sum_{v'} \exp(z_{v'})}
  && \text{($\exp(c) \neq 0$, cancel)} \\
  &= \mathrm{softmax}(z)_v.
\end{align}
\end{derivation}

\Cons\ Shift invariance justifies the \emph{log-sum-exp} (LSE) trick
\begin{equation}
\log q_\theta(x_t = v \mid x_{<t}) = z_v - \underbrace{\log \sum_{v'} \exp(z_{v'})}_{=:\,\mathrm{LSE}(z)}.
\end{equation}
Naively computing $\exp(z_v)$ for large logits overflows 32-bit floats;
subtracting $\max_v z_v$ before exponentiating is numerically equivalent
by invariance and keeps intermediate values bounded.

\Cons\ Softmax couples the probabilities of all tokens: increasing the
logit of one token necessarily reduces the probability of every other
token (the total must sum to one). A probability of $0.9$ on token $A$
does not mean ``the model is 90\% confident in $A$ as a fact''---it
means ``$A$ dominates all other tokens combined by a factor of roughly
$9{:}1$ under this parameterisation.'' The relative nature of softmax
probabilities is the source of many calibration problems we examine
below.

\subsection{Worked Example: Softmax and LSE}

\begin{workedexample}[: softmax and LSE on a toy vocabulary]
Take $V = \{\text{\emph{yes}}, \text{\emph{no}}, \text{\emph{maybe}}\}$
and logits $z = (2.0,\ 1.0,\ 0.1)$.
\begin{enumerate}
  \item Unnormalised exponentials: $(\exp 2,\ \exp 1,\ \exp 0.1)
  \approx (7.389,\ 2.718,\ 1.105)$, summing to $Z \approx 11.212$.
  \item Probabilities:
  $q(\text{\emph{yes}}) \approx 0.659$,
  $q(\text{\emph{no}}) \approx 0.242$,
  $q(\text{\emph{maybe}}) \approx 0.099$.
  \item If the ground-truth token is ``yes'', per-token NLL is
  $-\log 0.659 \approx 0.417\,\nats$.
  \item Shift every logit by $+1000$. By the derivation above the
  probabilities are numerically identical, but naive
  $\exp$-then-normalise overflows $\mathrm{float32}$. Subtracting the
  max ($+1002$ after shift) returns to shifted logits
  $(0,\ -1,\ -1.9)$, which is exactly $z - \max z$ in the original
  frame.
\end{enumerate}
Shift invariance is not a theoretical nicety; it is a daily
engineering tool.
\end{workedexample}

\subsection{The Decomposition of Cross-Entropy Into Per-Position Losses}

\Cons\ Training divides cleanly along positions. Because equation
(\ref{eq:mle}) sums over positions $t$ as well as over sequences, the
total loss on a batch decomposes exactly into a sum of per-position NLL
contributions, each equivalent to a multinomial logistic regression
problem over $V$. This decomposition is what makes teacher-forcing
training efficient: a single forward pass computes the conditional
distribution at every position in parallel, and backpropagation
distributes the gradient across all of them simultaneously.

\Emp\ A mis-tokenised example (one where a rare token has been split
unusually) contributes anomalously large loss spikes to specific
positions---a phenomenon we revisit in Chapter~4.

\begin{provedassumed}
\textbf{Proved here:} the equivalences
(\ref{eq:mle})$\Leftrightarrow$(\ref{eq:nll-ce}), the identity
(\ref{eq:ce-kl}) $H(p,q) = H(p) + \KL(p\parallel q)$, and the softmax
shift invariance of (\ref{eq:softmax}).
\textbf{Assumed here:} A1.1b (token-pair averaging), A1.2 (strictly positive
softmax), and the chain-rule factorisation (\ref{eq:chain}).
\textbf{Empirically observed (not proved):} that cross-entropy
training on large corpora yields models whose outputs \emph{look}
coherent and contextually appropriate; this is a property of the
data distribution, not a theorem.
\textbf{Used later:} Chapter~2 optimises (\ref{eq:mle}); Chapter~5
parameterises $q_\theta$ via attention; Chapter~9 replaces $\hat p$ with
a preference distribution.
\end{provedassumed}

\section{Information-Theoretic Foundations}

To reason about training curves, perplexity benchmarks, and calibration, we need a modest amount of classical information theory. The canonical reference is \citet{cover2006elements}; the textbook treatment is \citet{bishop2006prml} or \citet{murphy2022pml}.

\subsection{Entropy}

\Def\ Let $p$ be a distribution on a discrete set $\mathcal{X}$. The
\emph{entropy} of $p$ is
\begin{equation}
H(p) = -\sum_{x \in \mathcal{X}} p(x) \log p(x)
     = \E_{x \sim p}[-\log p(x)].
\end{equation}
With natural log the unit is $\nats$; with $\log_2$ the unit is $\bits$.
(We use $\nats$ throughout unless stated.)

\Ident\ Entropy is non-negative, equals $0$ iff $p$ is a point mass,
and is maximised on a finite support of size $|V|$ by the uniform
distribution, for which $H = \log |V|$.

\Cons\ Entropy measures the average surprise under $p$.
\citet{shannon1948} interpreted $H(p)$ in $\bits$ as the expected
number of bits required to encode a random draw from $p$ under the
optimal code (Shannon's source-coding theorem).

\subsection{Cross-Entropy and KL Divergence}

\Def\ The \emph{cross-entropy} between $p$ and $q$ (both on the same
finite support, with $q > 0$) is
\begin{equation}
H(p, q) = -\sum_{x} p(x) \log q(x) = \E_{x \sim p}[-\log q(x)],
\end{equation}
interpreted as the expected coding cost of samples from $p$ using a
code optimised for $q$.

\Def\ The \emph{Kullback--Leibler divergence} from $p$ to $q$ is
\begin{equation}
\KL(p \parallel q) = H(p, q) - H(p) = \sum_x p(x) \log \frac{p(x)}{q(x)}.
\label{eq:kl-def}
\end{equation}

\Ident\ $\KL$ is non-negative, zero iff $p = q$, and asymmetric (in
general $\KL(p\parallel q) \neq \KL(q\parallel p)$). Asymmetry reappears
in Chapter~9 where reverse-KL and forward-KL yield qualitatively different
alignment objectives.

\Cons\ The identity $H(p,q) = H(p) + \KL(p\parallel q)$ (already
derived as (\ref{eq:ce-kl})) expresses $H(p,q)$ as the sum of an
irreducible part $H(p)$ and a \emph{coding-overhead} part
$\KL(p\parallel q)$. Cross-entropy training minimises the overhead.

\subsection{Mutual Information}

\Def\ For random variables $X, Y$ with joint $p(x, y)$ and marginals
$p(x), p(y)$, the mutual information is
\begin{equation}
I(X; Y) := \KL\!\big(p(x, y) \parallel p(x)p(y)\big) = H(X) - H(X \mid Y).
\end{equation}

\Cons\ $I(X; Y)$ quantifies how much uncertainty about $X$ is resolved
by observing $Y$. In LLM work, mutual information appears in
information-bottleneck arguments about representation learning, in
the analysis of in-context learning (where the context tokens reduce
uncertainty about the continuation), and in retrieval evaluation
(where good retrieval increases $I(\text{answer}; \text{context})$).

\subsection{Perplexity}

\Def\ The \emph{perplexity} of $q_\theta$ on $\mathcal{D}$ is the
exponential of per-token cross-entropy:
\begin{equation}
\mathrm{PPL}(q_\theta; \mathcal{D})
  := \exp\!\left( \frac{1}{N_{\mathrm{tok}}} \sum_{i,t}
       -\log q_\theta\!\big(x^{(i)}_t \mid x^{(i)}_{<t}\big) \right)
  = \exp\!\big( H(\hat p,\, q_\theta) \big).
\label{eq:perplexity}
\end{equation}

\Ident\ Because $\exp : \mathbb{R} \to \mathbb{R}_{>0}$ is strictly
monotone,
\begin{equation}
\arg\min_\theta\, H(\hat p,\, q_\theta)
  \;=\; \arg\min_\theta\, \exp\!\big(H(\hat p,\, q_\theta)\big)
  \;=\; \arg\min_\theta\, \mathrm{PPL}(q_\theta; \mathcal{D}).
\label{eq:ppl-argmin}
\end{equation}
Minimising cross-entropy and minimising perplexity define the same
optimisation problem.

\Cons\ Perplexity admits a heuristic reading as an ``effective
vocabulary size'' the model is uncertain over at each step --- the
size of a uniform distribution that would have the same entropy.
This is intuition only, not a theorem: a perplexity of $20$ does
not imply the model is choosing among $20$ specific tokens. It does
imply the conditional distribution has the same entropy as the
uniform on $20$ symbols. A perplexity of $1$ corresponds to a
perfectly deterministic predictor; a perplexity of $|V|$ corresponds
to the uniform distribution over the full vocabulary. The
exponential scale makes small differences in cross-entropy (in $\nats$)
visually salient on a log-plot of PPL.

\Emp\ Perplexity is \emph{tokenization-dependent}: a model with a
smaller vocabulary mechanically achieves lower per-token perplexity
because each token is less informative. Comparing perplexities across
models trained with different tokenisers is therefore misleading; this
is revisited in Chapter~4.

\Emp\ A low perplexity does not imply high task performance. A model
can be excellent at modelling the statistics of Wikipedia prose while
still being useless at answering a question. Chapter~10 confronts task-level
evaluation directly.

\begin{provedassumed}
\textbf{Proved here:} the CE/KL identity (\ref{eq:ce-kl}) and the
cross-entropy/perplexity equivalence (\ref{eq:ppl-argmin}).
\textbf{Assumed here:} finite support, strictly positive $q$ (A1.2),
and $\log$ in $\nats$.
\textbf{Empirically observed:} perplexity is tokenization-dependent and
only weakly correlated with downstream task accuracy.
\textbf{Used later:} Chapter~6 uses (\ref{eq:perplexity}) as a headline
metric; Chapter~9 uses (\ref{eq:kl-def}) as the KL penalty in DPO.
\end{provedassumed}

\section{Uncertainty, Confidence, and Calibration}

\subsection{Three Quantities That Look the Same and Are Not}

\Def\ A common source of error in LLM engineering is the collapse of
three distinct quantities into a single intuitive notion of ``the model
knows.'' They are:
\begin{itemize}
  \item \textbf{Probability}: the scalar $q_\theta(v \mid x_{<t})$ the
  model assigns to a specific token $v$ given its parameters and
  context. This is the output of the softmax (\ref{eq:softmax}).
  \item \textbf{Confidence}: a summary statistic of the full
  distribution $q_\theta(\cdot \mid x_{<t})$, typically operationalised
  as $1 - H(q_\theta) / \log |V|$ (normalised entropy), or as $\max_v
  q_\theta(v \mid x_{<t})$ (top-1 probability), or as $\log q_\theta(v_1)
  - \log q_\theta(v_2)$ (log-odds between top-2).
  \item \textbf{Correctness}: whether the sampled or argmax token
  matches an external ground truth. Correctness is defined relative to
  something outside the model---a labeller, a gold answer, a runtime
  assertion.
\end{itemize}

\Emp\ These three can, and routinely do, move independently. A model
can assign $0.9$ probability to the correct token (high probability,
high confidence, correct); assign $0.9$ to an incorrect token (high
probability, high confidence, \emph{incorrect}); or spread probability
evenly over five plausible tokens of which one is correct (low
confidence, correct-on-average). The difference between the second and
third cases is the entire subject of \emph{calibration}.
Figure~\ref{fig:prob_conf_corr} illustrates the three cases side by
side.

\begin{figure}[H]
  \centering
  \begin{tikzpicture}[
  case/.style={draw, rounded corners, minimum width=4.2cm, minimum height=2.2cm,
               align=left, font=\scriptsize},
  bar/.style={draw, thick},
  ok/.style={fill=green!30},
  bad/.style={fill=red!30},
  node distance=0.35cm
]

% Case A: high prob, high confidence, correct
\node[case, fill=green!6] (A) {
  \textbf{Case A} (probability, confidence, correctness all aligned)\\[2pt]
  \textit{Top-1} ``Paris'': $0.92$\\
  \textit{next} ``Rome'': $0.03$\\
  \textit{next} ``Berlin'': $0.02$\\[2pt]
  $\Rightarrow$ correct \& well-calibrated
};

% Case B: high prob, high confidence, wrong
\node[case, fill=red!6, right=of A] (B) {
  \textbf{Case B} (high confidence, \emph{wrong})\\[2pt]
  \textit{Top-1} ``Sydney'': $0.88$\\
  \textit{next} ``Canberra'': $0.05$\\
  \textit{next} ``Melbourne'': $0.04$\\[2pt]
  $\Rightarrow$ hallucination (overconfident)
};

% Case C: low confidence, correct-on-average
\node[case, fill=yellow!10, right=of B] (C) {
  \textbf{Case C} (low confidence, correct-on-avg)\\[2pt]
  \textit{Top-1} ``2013'': $0.21$\\
  \textit{next} ``2012'': $0.20$\\
  \textit{next} ``2014'': $0.18$\\[2pt]
  $\Rightarrow$ calibrated uncertainty
};

% Three-axis legend below
\node[font=\footnotesize\bfseries, below=0.6cm of B.south] (legend)
  {Three orthogonal quantities:};
\node[font=\footnotesize, below=0.05cm of legend]
  {\textbf{probability} $q_\theta(v\mid x_{<t})$ $\cdot$
   \textbf{confidence} $\max_v q_\theta$ or $1 - H/\log|V|$ $\cdot$
   \textbf{correctness} $\mathbb{1}[\hat y = y]$};

\end{tikzpicture}

  \caption{Probability, confidence, and correctness are orthogonal
           axes. The vertical axis is the scalar
           $q_\theta(v^\star \mid x_{<t})$ from (\ref{eq:softmax}); the
           horizontal axis is the dispersion summary
           $1 - H(q_\theta)/\log|V|$ derived from the softmax output;
           the marker shape encodes correctness against external
           ground truth. Case A is aligned; Case B is the classic
           hallucination (high confidence, wrong); Case C is
           calibrated uncertainty in the presence of genuine
           ambiguity.}
  \label{fig:prob_conf_corr}
\end{figure}

\subsection{Calibration and Reliability Diagrams}

\Def\ A probabilistic classifier is \emph{perfectly calibrated} if,
for every $p^\star \in [0, 1]$,
\begin{equation}
\Pr\!\big[\,\text{correct} \,\big|\, \hat q = p^\star\,\big] = p^\star,
\label{eq:perfect-calibration}
\end{equation}
where $\hat q$ is the model's top-1 probability on the prediction and
``correct'' is the indicator that the top-1 token matches the label.

\Ident\ Perfect calibration is an \emph{extra} property beyond
accuracy: a classifier can be accurate but uncalibrated (its
confidences misstate reliability) or calibrated but inaccurate
(uniformly low confidence on a hard problem).

\Def\ A \emph{reliability diagram} is the empirical estimator of
(\ref{eq:perfect-calibration}). Partition $[0,1]$ into $K$ bins
$\mathcal{I}_1, \ldots, \mathcal{I}_K$ and define, for each bin $k$,
\begin{align}
B_k &:= \{\, i : \hat q_i \in \mathcal{I}_k \,\}, &
\mathrm{acc}(B_k) &:= \frac{1}{|B_k|} \sum_{i \in B_k} \mathbb{1}[\hat y_i = y_i], &
\mathrm{conf}(B_k) &:= \frac{1}{|B_k|} \sum_{i \in B_k} \hat q_i.
\label{eq:bin-defs}
\end{align}
Plotting $\mathrm{acc}(B_k)$ against $\mathrm{conf}(B_k)$ for $k = 1,
\ldots, K$ gives the diagram in Figure~\ref{fig:reliability}. Points
\emph{below} the diagonal (predicted $>$ actual) indicate
overconfidence; points above indicate underconfidence.

\Emp\ Modern deep networks, including LLMs, are almost invariably
overconfident in the high-probability regime
\citep{guo2017calibration}.

\begin{figure}[H]
  \centering
  \begin{tikzpicture}[
  scale=1.0,
  axis/.style={->, thick, >=Stealth},
  bar/.style={fill=blue!30, draw=blue!60},
  calline/.style={thick, black, dashed},
  overconf/.style={thick, red, smooth},
  underconf/.style={thick, teal, smooth},
  lbl/.style={font=\footnotesize}
]

% Axes
\draw[axis] (0,0) -- (6.2,0) node[right, lbl] {confidence $\operatorname{conf}(B_k)$};
\draw[axis] (0,0) -- (0,6.2) node[above, lbl] {accuracy $\operatorname{acc}(B_k)$};

% Tick marks + labels
\foreach \x/\t in {0/0, 3/0.5, 6/1} {
  \draw (\x,0) -- (\x,-0.12) node[below, lbl] {\t};
  \draw (0,\x) -- (-0.12,\x) node[left, lbl] {\t};
}

% Perfect calibration diagonal
\draw[calline] (0,0) -- (6,6) node[pos=0.7, above left, lbl, black] {perfect calibration};

% Overconfident curve (below diagonal): acc < conf in high-confidence region
\draw[overconf] plot[domain=0:6, samples=60]
  (\x, {\x - 0.07*\x*\x - 0.003*\x*\x*\x});
\node[lbl, red, right, align=left] at (5.0, 2.0) {overconfident\\(modern DNNs)};

% Underconfident curve (above diagonal): acc > conf
\draw[underconf] plot[domain=0:6, samples=60]
  (\x, {\x + 0.07*\x*(6-\x)});
\node[lbl, teal, above] at (2.2, 4.2) {underconfident};

% Histogram bars at the bottom showing predicted-prob distribution
\foreach \i/\h in {0.5/0.12, 1.1/0.18, 1.7/0.22, 2.3/0.22, 2.9/0.25, 3.5/0.32, 4.1/0.40, 4.7/0.55, 5.3/0.78, 5.9/0.95}{
  \pgfmathsetmacro\hh{\h*0.5}
  \draw[bar] (\i-0.25,-0.55) rectangle (\i+0.25,-0.55+\hh);
}
\node[lbl, below right=0.6cm and -0.3cm of current bounding box.south west] {\textit{density of predictions at each confidence level (toy)}};

\end{tikzpicture}

  \caption{Reliability diagram. The dashed black diagonal is the
           locus (\ref{eq:perfect-calibration}) of perfect calibration.
           The red curve is the overconfident regime
           ($\mathrm{acc}(B_k) < \mathrm{conf}(B_k)$), typical for
           modern deep networks; the teal curve is underconfident
           ($\mathrm{acc}(B_k) > \mathrm{conf}(B_k)$). The vertical
           bars below each bin show the bin count $|B_k|$; bins with
           few samples (small $|B_k|$) have high-variance accuracy
           estimates and should not be over-interpreted.}
  \label{fig:reliability}
\end{figure}

\subsection{Expected Calibration Error}

\Def\ Let the \emph{calibration function} $c(p) := \Pr[\text{correct}
\mid \hat q = p]$ denote the conditional accuracy of the predictor
when its top-class probability is $p$ (defined a.e.\ with respect to
the distribution of $\hat q$). The \emph{Expected Calibration Error}
(ECE) is the $L_1$-weighted gap between this calibration function and
the diagonal:
\begin{equation}
\mathrm{ECE}(q; \mathcal{D}_{\mathrm{cal}})
  := \E_{\hat q}\!\big[\,\big| c(\hat q) - \hat q \big|\,\big].
\label{eq:ece-population}
\end{equation}
Because $\hat q$ is continuous in general, the conditioning
$\Pr[\text{correct} \mid \hat q = p]$ should be read as the
regular-conditional version, and the operational estimator below
replaces the point conditioning by binning into neighbourhoods
$\mathcal{I}_k$.

\begin{derivation}[: finite-bin ECE estimator]
Replace the expectation in (\ref{eq:ece-population}) by its empirical
version over a calibration split $\mathcal{D}_{\mathrm{cal}}$ of size
$N$, using the binning (\ref{eq:bin-defs}):
\begin{align}
\widehat{\mathrm{ECE}}_{K}
  &= \sum_{k=1}^{K} \widehat{\Pr}[\,\hat q \in \mathcal{I}_k\,]
       \cdot \big| \widehat{\Pr}[\,\text{correct} \mid \hat q \in \mathcal{I}_k\,] - \E[\hat q \mid \hat q \in \mathcal{I}_k] \big| \\
  &= \sum_{k=1}^{K} \frac{|B_k|}{N}\,
       \big| \mathrm{acc}(B_k) - \mathrm{conf}(B_k) \big|.
\label{eq:ece}
\end{align}
The replacement uses $\widehat{\Pr}[\hat q \in \mathcal{I}_k] = |B_k|/N$
and plug-in estimators for the conditional expectations. Empty bins
($|B_k| = 0$) contribute zero to the sum by the leading
$|B_k|/N$ weight; their $\mathrm{acc}$ and $\mathrm{conf}$ are
left undefined. (Equivalently, restrict the sum to bins with
$|B_k| > 0$.)
\end{derivation}

\Cons\ A well-calibrated predictor has
$\widehat{\mathrm{ECE}}_K \approx 0$; values above a few percent
indicate practically significant miscalibration. ECE is commonly
reported with $K = 10$ or $K = 15$ equal-width bins.

\paragraph{Small-sample caveat.}
\Assum\ The estimator (\ref{eq:ece}) assumes each bin contains enough
samples that $\mathrm{acc}(B_k)$ is a good estimate of the underlying
conditional accuracy. As a rule of thumb, bins with $|B_k| \lesssim 30$
have standard errors on $\mathrm{acc}(B_k)$ of order
$1/\sqrt{|B_k|} \gtrsim 0.18$, which can drive $\widehat{\mathrm{ECE}}_K$
entirely by binning noise. Equal-\emph{mass} binning (quantile-based
$\mathcal{I}_k$) mitigates this at the cost of bin locations that shift
between evaluations. \citet{kumar2019ece} and \citet{nixon2019measuring}
discuss bias and variance of (\ref{eq:ece}) in detail.

\paragraph{End-to-end ECE: a $K=3$, $N=12$ walk-through with an
empty bin.}
\begin{workedexample}[: bin-by-bin computation including an empty bin]
Fix $K = 3$ equal-width bins
$\mathcal{I}_1 = [0,\,1/3]$, $\mathcal{I}_2 = (1/3,\,2/3]$,
$\mathcal{I}_3 = (2/3,\,1]$ over $N=12$ predictions. Suppose the
predicted top-class probabilities $\hat q_i$ and the indicators
$\mathbb{1}[\hat y_i = y_i]$ fall as follows:
\begin{center}
\small
\begin{tabular}{|c|l|c|}
\hline
Bin & $\hat q_i$ values & correct? \\
\hline
$B_1$ ($|B_1|=4$) & $\{0.10, 0.15, 0.25, 0.30\}$ & $\{0,1,1,0\}$ \\
$B_2$ ($|B_2|=0$) & --- & --- \\
$B_3$ ($|B_3|=8$) & $\{0.75, 0.80, 0.82, 0.85, 0.88, 0.90, 0.92, 0.95\}$ & $\{1,1,1,0,1,1,1,1\}$ \\
\hline
\end{tabular}
\end{center}
Compute per-bin accuracy and confidence:
\begin{align*}
\mathrm{acc}(B_1) &= 2/4 = 0.50, &
\mathrm{conf}(B_1) &= (0.10+0.15+0.25+0.30)/4 = 0.20, \\
\mathrm{acc}(B_3) &= 7/8 = 0.875, &
\mathrm{conf}(B_3) &= (0.75+0.80+0.82+0.85+0.88+0.90+0.92+0.95)/8 \approx 0.859.
\end{align*}
$B_2$ is empty: $|B_2| = 0$ contributes zero to the sum by the
leading $|B_k|/N$ weight, regardless of $\mathrm{acc}$ and
$\mathrm{conf}$ (which are undefined there). The estimator is then
\[
\widehat{\mathrm{ECE}}_3
  \;=\; \tfrac{4}{12}\,|0.50 - 0.20|
        + \tfrac{0}{12}\,(\text{undefined})
        + \tfrac{8}{12}\,|0.875 - 0.859|
  \;\approx\; 0.100 + 0 + 0.011
  \;=\; 0.111.
\]
The walk-through illustrates three details the formal expression
hides: that empty bins drop out automatically, that bin~1 is
underconfident (accuracy $>$ confidence) while bin~3 is mildly
overconfident, and that the small-sample caveat above applies to
$B_1$ ($|B_1|=4 \ll 30$): its accuracy estimate has standard error
of order $0.25$, so the $0.10$ contribution to $\widehat{\mathrm{ECE}}_3$
is itself uncertain at the $\pm 0.08$ level.
\end{workedexample}

\paragraph{ECE fails to discriminate risk: a two-bin counterexample.}
\begin{workedexample}[: same ECE, very different risk]
Fix $K = 2$ bins $\mathcal{I}_1 = [0, 0.5]$ and
$\mathcal{I}_2 = (0.5, 1]$ and take $N = 100$ samples.

\textbf{System A.} 50 predictions in each bin. In bin~1,
$\mathrm{conf} = 0.30$, $\mathrm{acc} = 0.50$ (underconfident by $0.20$).
In bin~2, $\mathrm{conf} = 0.90$, $\mathrm{acc} = 0.70$ (overconfident
by $0.20$). Then
\[
\widehat{\mathrm{ECE}}_2(A) = \tfrac{50}{100} \cdot 0.20 + \tfrac{50}{100} \cdot 0.20 = 0.20.
\]

\textbf{System B.} 50 predictions in each bin. In bin~1,
$\mathrm{conf} = 0.30$, $\mathrm{acc} = 0.50$ (underconfident by
$0.20$). In bin~2, $\mathrm{conf} = 0.70$, $\mathrm{acc} = 0.90$
(also underconfident by $0.20$). Then
\[
\widehat{\mathrm{ECE}}_2(B) = \tfrac{50}{100} \cdot 0.20 + \tfrac{50}{100} \cdot 0.20 = 0.20.
\]

\textbf{Interpretation.} Both systems have identical ECE. Yet an
operating rule ``abstain unless $\hat q > 0.5$'' yields error rate
$0.30$ on system~A (accepts bin~2, which is only $0.70$ accurate) and
error rate $0.10$ on system~B (accepts bin~2, which is $0.90$
accurate). ECE is signed-insensitive: it weights
$|\mathrm{acc} - \mathrm{conf}|$ without regard to direction, even
though overconfidence and underconfidence have opposite consequences
for any threshold-on-probability decision rule.
\end{workedexample}

\Cons\ ECE has known weaknesses: sensitivity to binning
\citep{kumar2019ece, nixon2019measuring}, insensitivity to the sign of
miscalibration (as above), and conflation of different failure modes.
Alternatives include the \emph{Brier score}
\begin{equation}
\mathrm{Brier}(q; \mathcal{D})
  := \frac{1}{N} \sum_{i=1}^{N} \sum_{v \in V}
       \big( q_i(v) - \mathbb{1}[y_i = v] \big)^2,
\label{eq:brier}
\end{equation}
and the marginal log-loss (per-instance NLL). Both are proper scoring
rules in the sense of the next subsection.

\subsection{Temperature Scaling}

\Def\ Let $\tau > 0$ be a scalar. \emph{Temperature scaling} divides
logits by $\tau$ before the softmax:
\begin{equation}
q_\tau(v \mid x_{<t}) = \frac{\exp(z_v / \tau)}{\sum_{v'} \exp(z_{v'} / \tau)}.
\label{eq:temperature}
\end{equation}
Temperature $\tau < 1$ sharpens the distribution (more mass on the
top-1); $\tau > 1$ flattens it; $\tau \to 0^+$ degenerates to a point
mass on the argmax; $\tau \to \infty$ degenerates to the uniform
distribution.

\Ident\ Temperature preserves the \emph{ordering} of all tokens: the
map $z \mapsto z/\tau$ is strictly monotone, so
$\arg\max q_\tau = \arg\max q_1$ for every $\tau > 0$. Temperature
cannot turn a low-probability token into the most probable one.

\Emp\ Fitting a single scalar $\tau$ on held-out data to minimise NLL
recovers most of the calibration deficit in image classifiers and
language models \citep{guo2017calibration}. Because $\tau$ is scalar,
argmax predictions are unchanged: calibration improves without
changing accuracy.

\Cons\ Temperature is also the primary decoding-time knob in
inference stacks: the same operation that fixes calibration can be
used to trade off diversity for determinism in generation. Temperature
does not introduce new knowledge; it redistributes existing mass.

\subsection{Proper Scoring Rules}

\Def\ A \emph{scoring rule} $S(q, y)$ assigns a score to a predicted
distribution $q$ in light of an observed outcome $y$. A scoring rule is
\emph{proper} if reporting the true distribution is, in expectation,
optimal:
\begin{equation}
\E_{y \sim p}[S(p, y)] \le \E_{y \sim p}[S(q, y)] \quad \text{for all } q.
\end{equation}
(The inequality is with respect to losses; reverse it if $S$ is a
gain.)

\Ident\ Log-loss $S(q, y) = -\log q(y)$ and the Brier score
(\ref{eq:brier}) are both proper. The canonical survey is
\citet{gneiting2007strictly}.

\Cons\ This is the theoretical backstop for cross-entropy training:
we are not just minimising an arbitrary convex surrogate; we are
minimising a proper scoring rule that is optimal \emph{in expectation}
when $q = p$. Equivalently, the uniquely optimal strategy under
log-loss is to report the true conditional distribution. The gap
between that optimum and what a trained model actually does tells us
how far we are from the optimum---and, informally, where there is room
to grow.

\begin{provedassumed}
\textbf{Proved here:} the finite-bin ECE estimator (\ref{eq:ece}) as
the plug-in of (\ref{eq:ece-population}); the construction of two
systems with identical ECE but opposite sign of miscalibration.
\textbf{Assumed here:} sufficient bin counts $|B_k|$ for the plug-in
to be a good estimator (A1.3 no distribution shift between
calibration and operating distributions).
\textbf{Empirically observed:} modern deep networks are overconfident
in the high-probability regime; temperature scaling recovers most of
the calibration deficit.
\textbf{Used later:} Chapter~6 revisits temperature as a decoding knob;
Chapter~9 tracks calibration drift after alignment; Chapter~10 uses
calibration as part of abstention and tool-use policies.
\end{provedassumed}

\section{Sampling and Decoding as Probabilistic Objects}

\Def\ Once the model has computed $q_\theta(\cdot \mid x_{<t})$, a
\emph{decoding strategy} selects a next token. Formally, a decoding
strategy is a (possibly stochastic) map from $q_\theta(\cdot \mid
x_{<t})$ to a token in $V$. This step is frequently treated as
incidental, but it materially shapes system behaviour.

\begin{itemize}
  \item \textbf{Greedy decoding} selects $\arg\max_v q_\theta(v \mid
  x_{<t})$. It is deterministic but pathological in long sequences, as
  it cannot recover from early mistakes and tends to produce repetitive,
  degenerate text \citep{holtzman2020curious}.
  \item \textbf{Stochastic sampling} draws $x_t \sim q_\theta(\cdot
  \mid x_{<t})$. This yields diverse outputs but can produce
  low-probability (and low-quality) tokens in the long tail.
  \item \textbf{Top-$k$ sampling} restricts sampling to the $k$
  highest-probability tokens, renormalises, and samples. It cuts the
  tail but ignores the shape of the head.
  \item \textbf{Top-$p$ (nucleus) sampling} restricts to the smallest
  set whose cumulative mass exceeds a threshold $p$, and samples from
  that set. It adapts to the shape of the distribution.
  \item \textbf{Temperature} modifies any of the above by pre-scaling
  the logits. $\tau \to 0^+$ recovers greedy; $\tau \to \infty$ flattens
  to uniform.
\end{itemize}

\Cons\ Decoding strategies define a \emph{second distribution} over
outputs, layered on top of the model's learned conditional. The
question ``what does the model say?'' is under-specified without a
decoding strategy. Chapter~6 treats decoding more formally; here the
point is that sampling is a first-class design choice with measurable
effects on calibration, diversity, and downstream task performance.

\section{Common Probabilistic Failure Modes}

With the probabilistic framing in place, several characteristic LLM failure modes resolve into familiar statistical pathologies.

\paragraph{Confident fabrication (hallucination).}
\Emp\ The model assigns high probability to a token sequence that is
factually wrong. Under the probabilistic lens this is unremarkable: the
training objective (\ref{eq:mle}) rewards high probability on
\emph{corpus-likely} sequences, and a fluent fabrication can be
corpus-likely if the corpus contains many fluent fabrications in
adjacent contexts. No component of cross-entropy training pushes the
model toward truth; truth is not a variable in the optimisation.

\paragraph{Overconfidence in rare events.}
\Emp\ Modern deep networks tend to concentrate mass aggressively. Argmax
probabilities drift toward 1 as models are scaled and trained longer,
beyond what the empirical accuracy in the corresponding bin justifies.
This is the phenomenon that reliability diagrams expose
(Figure~\ref{fig:reliability}) and that temperature scaling partially
corrects.

\paragraph{Token-boundary brittleness.}
\Emp\ Because cross-entropy is computed per token and tokenisation is
lossy (Chapter~4), semantically equivalent inputs can produce very
different logit trajectories. The model may be confident on one phrasing
of a prompt and unsure on a paraphrase. This is a calibration failure
expressed at the level of inputs rather than outputs.

\paragraph{Exposure bias.}
\Cons\ During training the model conditions on ground-truth prefixes
(teacher forcing). At inference it conditions on its own generations.
Any probability mass on wrong tokens at position $t$ compounds over
positions $t+1, t+2, \ldots$, producing distributional drift that the
loss did not penalise. This is a structural consequence of the
decomposition (\ref{eq:mle}) and a reason why long-form generation is
harder than short-form.

\paragraph{Calibration drift after alignment.}
\Emp\ Post-training techniques (instruction tuning, RLHF; Chapter~9)
reshape \emph{which} tokens receive probability mass without preserving
\emph{how much}. Aligned models are often less calibrated than their
base counterparts; the confidence signal that was useful in the base
model becomes misleading after alignment. This is a critical
consideration in production pipelines and motivates retrieval-augmented
and tool-using systems (Chapter~10).

\begin{workedexample}[: overconfidence vignette]
A factual-QA benchmark reports an LLM's top-1 predictions at $88\%$
average top-1 probability but only $71\%$ top-1 accuracy. The
$17$-percentage-point gap is the calibration error. A downstream
system that uses top-1 probability as a confidence-for-abstention
signal---abstaining whenever probability is below $0.7$---never
abstains on the hard cases, because probability clusters above $0.7$
even when accuracy is poor. Fixing this is primarily a calibration
task (temperature scaling on a held-out split), not a model-capacity
task.
\end{workedexample}

\Cons\ A useful system-design principle follows: \emph{if a downstream
decision depends on a threshold on model probability, validate that the
probability is calibrated on the operating distribution before
deployment.}

\section{Systems Engineering Implications}

From a systems engineering standpoint, an LLM should be treated as a stochastic subsystem with an opaque internal state and non-stationary performance. The following table summarizes the consequences of the statistical framing for each phase of a standard V-model:

\begin{center}
\begin{tabular}{|L{3.4cm}|L{8.0cm}|}
\hline
\textbf{Phase / Concern} & \textbf{Consequence of the Statistical Framing} \\
\hline
Requirements & Specifications must admit probabilistic outputs; ``the system shall return the correct answer'' is not enforceable against a probabilistic model. \\
\hline
Architecture & Downstream components must tolerate variance. Verification layers, retrieval, tool use, and human-in-the-loop review are architectural patterns motivated directly by the probabilistic substrate. \\
\hline
Verification \& Validation & Evaluation must report distributions, not point estimates. Calibration metrics (ECE, Brier) belong in V\&V alongside accuracy. \\
\hline
Operations & Distribution shift is inevitable: input distributions change over time. Continuous monitoring of loss, entropy, and acceptance rates is the operational analogue of calibration. \\
\hline
Incident Response & Failures are rarely deterministic bugs. Root-cause analysis must consider rare-event tail behavior and distributional assumptions. \\
\hline
\end{tabular}
\end{center}

These properties motivate the architectural patterns---verification layers, retrieval augmentation, structured output formats, and human-in-the-loop review---that Weeks 8 through 10 will develop in detail.

\section{Human Factors and Failure Modes}

A common early failure mode in LLM systems is \emph{automation bias}:
\begin{quote}
``The model sounds confident, therefore it must be correct.''
\end{quote}
This arises from a mismatch between human epistemic intuitions---where confident assertions come from knowledgeable interlocutors---and probabilistic text generation, where confidence is a statistical property of the next-token distribution. Humans infer reliability from tone; models produce tone as a by-product of fitting tone-heavy text. The fluency heuristic is a cousin: the more coherent the output reads, the more we trust it, independent of whether it is correct.

Effective first-line mitigations include structured outputs with explicit confidence fields, uncertainty visualization in user interfaces, and source attribution (retrieval citations, tool-call traces). We return to these in Chapter~10 in the context of system design and in Chapter~11 in the context of assurance cases.

\section{Bridges to Later Chapters}

The statistical machinery of this chapter is revisited, refined, or challenged in every subsequent chapter. Explicit bridges worth keeping in mind:
\begin{itemize}
  \item \textbf{Chapter 2 (Optimization)} studies how we actually find $\hat\theta$ in practice, and why the optimization dynamics of cross-entropy over deep networks defy the intuition of classical statistics.
  \item \textbf{Chapter 4 (Tokenization \& Embeddings)} makes precise what $x_t$ is as an object: a token rather than a word, and embedded in a vector space rather than a categorical index.
  \item \textbf{Chapter 5 (Transformers)} specifies how the conditional $q_\theta(\cdot \mid x_{<t})$ is computed from the prefix, filling in the black box that the derivations above took for granted.
  \item \textbf{Chapter 6 (Pretraining \& ICL)} revisits perplexity, decoding, and the emergence of in-context adaptation as a property of the pretraining objective.
  \item \textbf{Chapter 8 (Adaptation)} and \textbf{Chapter 9 (Alignment)} change $q_\theta$ after pretraining. Both typically \emph{reduce} calibration relative to the base model, which forces the discussion from this chapter to be revisited operationally.
  \item \textbf{Chapter 10 (RAG \& Tools)} substitutes $q_\theta(y \mid x)$ with $q_\theta(y \mid x, d)$ where $d$ is retrieved evidence, converting a pure language-modelling task into a conditional inference task with grounded context.
  \item \textbf{Chapter 11 (Governance \& Assurance)} treats probabilistic performance as an assurance-case claim and asks what evidence is needed to justify it.
\end{itemize}

\section{Key References}

\begin{itemize}
  \item \citet{bishop2006prml} --- \emph{Pattern Recognition and Machine Learning}; definitive treatment of probabilistic modelling and calibration.
  \item \citet{murphy2012mlpp} and \citet{murphy2022pml} --- \emph{Machine Learning: A Probabilistic Perspective} and the updated \emph{Probabilistic Machine Learning}; covers MLE, exponential families, and information theory.
  \item \citet{shannon1948} --- \emph{A Mathematical Theory of Communication}; the foundational paper introducing entropy, cross-entropy, and the source-coding theorems that underpin perplexity.
  \item \citet{cover2006elements} --- \emph{Elements of Information Theory}; the canonical reference for information-theoretic quantities.
  \item \citet{vapnik1998statistical} --- \emph{Statistical Learning Theory}; classical PAC framework for generalization, useful background for Chapter~2.
  \item \citet{guo2017calibration} --- empirical study of calibration in modern neural networks; introduces temperature scaling.
  \item \citet{gneiting2007strictly} --- survey of proper scoring rules in forecasting.
  \item \citet{holtzman2020curious} --- shows that maximum-likelihood decoding fails in open-ended generation; motivates the decoding toolkit.
\end{itemize}

\section{Recommended University Lectures}

\begin{itemize}
  \item Caltech CS156 --- \emph{Learning From Data} \citep{caltech_cs156}.
  \item MIT 6.036 --- \emph{Introduction to Machine Learning} \citep{mit_6036}.
  \item Stanford CS229 --- \emph{Machine Learning}, particularly the probability and optimization sections \citep{stanford_cs229}.
\end{itemize}

\section{Practitioner Lab}

\textbf{Objective}: implement multinomial logistic regression for a toy next-token task and observe probabilistic behavior first-hand.

\textbf{Tasks}:
\begin{enumerate}
  \item Train a simple next-token classifier on a short corpus.
  \item Inspect the output probability distributions on held-out inputs; plot entropies and top-$k$ gaps.
  \item Measure accuracy, entropy, and calibration error (ECE) on a held-out split.
  \item Fit a scalar temperature on a validation split; re-measure ECE; verify that accuracy is unchanged.
  \item Sample generations at temperatures $\tau \in \{0.2, 0.7, 1.0, 1.4\}$ and qualitatively compare.
  \item Construct a case where top-1 probability exceeds $0.9$ but the top-1 prediction is wrong; discuss what a downstream system should do in response.
\end{enumerate}

\textbf{Deliverable}: a Jupyter notebook with code and plots, plus a one-page reflection titled ``Where my probabilistic intuition failed.''

See \texttt{labs/week01/week01\_lab\_probabilistic\_foundations.ipynb}.

\section*{Derivation Ledger (Ch. 1)}

\begin{derivledger}
A compact index of the key identities established in this chapter.
Each entry cites the full derivation already given in the text; no
new content.
\begin{itemize}
  \item \textbf{Chain-rule factorisation of sequences.}\
        $p_\theta(x_{1:T}) = \prod_{t=1}^{T} p_\theta(x_t \mid x_{<t})$
        --- Eq.~\eqref{eq:chain}.
  \item \textbf{MLE $\to$ token-normalised NLL $\to$ cross-entropy.}\
        Five-line derivation Eqs.~\eqref{eq:mle-1}--\eqref{eq:mle};
        assumption \textbf{(A1.1a)} i.i.d.\ sequences applied at line
        Eq.~\eqref{eq:mle-1}. Final form connects training loss to
        the cross-entropy Eq.~\eqref{eq:nll-ce}.
  \item \textbf{Cross-entropy identity.}\
        $H(p,q) = H(p) + \KL(p \,\|\, q)$
        --- Eq.~\eqref{eq:ce-kl}. Non-negativity of KL gives
        $H(p,q) \ge H(p)$, with equality iff $q = p$ a.e.
  \item \textbf{Softmax and shift invariance.}\
        $\mathrm{softmax}(z) = \mathrm{softmax}(z + c \mathbf{1})$
        --- Eq.~\eqref{eq:softmax}; $\exp(c)$ cancels in numerator
        and denominator.
  \item \textbf{Perplexity and ordering.}\
        $\mathrm{PPL} = \exp(H)$, Eqs.~\eqref{eq:perplexity}--\eqref{eq:ppl-argmin};
        monotone $\exp$ implies $\arg\min H = \arg\min \mathrm{PPL}$.
  \item \textbf{ECE (finite-bin estimator).}\
        $\mathrm{ECE} = \sum_{k=1}^{K} (|B_k|/N)\,|\mathrm{acc}(B_k) - \mathrm{conf}(B_k)|$
        --- Eq.~\eqref{eq:ece}; population form Eq.~\eqref{eq:ece-population}.
        Small-sample caveat carried with the estimator.
  \item \textbf{Brier score.}\
        Proper scoring-rule decomposition --- Eq.~\eqref{eq:brier}.
  \item \textbf{Temperature rescaling.}\
        Divides logits by $\tau$ before softmax --- Eq.~\eqref{eq:temperature}.
\end{itemize}
\end{derivledger}

\section{Chapter 1 Takeaway}

\begin{quote}
An LLM is not a knowledge base or a reasoner; it is a probabilistic sequence model trained by minimizing cross-entropy against a corpus distribution. Cross-entropy is a proper scoring rule, which makes it a principled training target---but it measures distributional fidelity, not truth. Everything downstream, from hallucination to overconfidence to the surprising power of prompting, is a consequence of reading the model's outputs as probabilities and not as assertions.
\end{quote}

All subsequent components---Transformers, scaling, adaptation, alignment, retrieval, tools, agents, and governance---build on this foundation.
